\begin{helper}
For every two-bite sample \(\omega\), the red part of
\(\noderef{TwoBiteFinalGraph}\) contains no triangle.
\end{helper}
\begin{proof}
Write \(G_R\) for the final red graph and \(\widetilde G_R\) for the red
overlay graph.  Suppose, for contradiction, that \(u,v,w\) span a triangle in
\(G_R\).  Since a final red edge is, by definition, an overlay red edge that was
not red-deleted, the three pairs \(uv,uw,vw\) are edges of \(\widetilde G_R\).

By \(\noderef{TwoBiteRedTriangleLatestEdge}\), applied to the three red overlay
edges \(uv,uw,vw\), one of the three red projected pairs is strictly latest in
the pair-rank order used by \(\noderef{TwoBiteFinalGraph}\).

If the largest rank is that of \(\{r_u,r_v\}\), then the red deletion predicate
for the directed edge \(uv\) holds with witness \(w\): the three overlay red
edges are \(uv,uw,vw\), and the two other projected-pair ranks are strictly
smaller than the rank of \(uv\).  Hence the undirected final-red deletion
predicate for \(uv\) holds, contradicting that \(uv\) remains an edge of
\(G_R\).  The cases in which the largest projected pair is \(\{r_u,r_w\}\) or
\(\{r_v,r_w\}\) are identical after permuting the names of the triangle
vertices.  In every case one of the three assumed final red edges is selected
by the red triangle deletion rule, a contradiction.
\end{proof}
